\bookchapter{第三十二章：手機安全急救手冊}

\chapterintro{此章乃一份「應急預案」。望爾永遠用不上，然果遇手機丟失、被盜、被詐騙之情事，此手冊能助爾將損失降至最低；地震、暴雨等災害來臨時，手機同為接收預警、定位求救、保持聯絡之工具。建議爾今即通讀一遍，然後將關鍵步驟截圖保存，或打印出來置於家中——蓋出事之時爾之手機或已不於手上，災時則可能沒電或沒信號。}

\booksection{場景零：地震／極端天氣——手機何以助爾保命}

所防者人身之安，非帳號資金。預警設置詳見第十三章「天氣預警、自然災害預警與地震預警」；此處只講警報響起之後該何以做。

\textbf{地震警報響了：先躲後撤，勿只顧拍視頻}\par

\begin{enumerate}
  \item 就近避險：趴下、掩護、抓牢（桌下、承重牆角、堅固傢具旁），護住頭頸；遠離玻璃、外牆、吊燈。
  \item 勿盲目奔樓：主震往往只有十幾秒至幾十秒，樓梯間可能更危險；電梯絕對不能用。
  \item 震停後再有序撤離：走樓梯、留意餘震與掉落物；到開闊地，遠離電線桿、廣告牌、危房。
  \item 手機用於緊要處：予緊急聯絡人發簡短消息（「我於XX，安全／受傷，須幫助」）較打長電話更省電；能發短信／微信時優先文字。
  \item 沒信號時：見第七章衛星SOS／衛星消息（機型支持時）；或到高處、開闊地再試；部分場景下無網藍牙對講亦有助。
  \item 關閉非必要之功：應急時啟飛行模式，定時開關；罷藍牙搜尋，降其亮度，可多支數時；留末格之電，以報平安與救援之訊。
\end{enumerate}

\begin{tcbwarn}{暴雨 / 颱風 / 山洪：預警之期，即當行動}
\begin{itemize}
  \item 見紅色預警：宜少外出，遠離地下之室、下穿之橋、河道、陡坡。
  \item 困於戶外：登堅固屋宇之高處，勿強涉水（積水或失井蓋、或漏電）；呼 119/110 時，詳陳地標與水位。
  \item 未行之先：以天氣 App 與當地預警之公號核其險；長途自駕，避山區臨水臨崖之路。
  \item 家人相報：約曰「若失聯，以最後微信、短信之定位為準」，免各自狂撥而耗盡其電。
\end{itemize}

\end{tcbwarn}

\begin{tcbpractice}{平素即須做之四事}
\begin{enumerate}
  \item 啟系統地震預警 / 正規預警 App，許其破靜音而出聲。
  \item 於通訊錄置頂二三人為「緊急聯絡人」，鎖屏之界面可顯之（iOS、安卓之醫療急救卡或鎖屏之訊）。
  \item 家中行囊，常備充電寶；遠行之前，滿其電而攜其線。
  \item 默識至少一號（家人之手機號），雖無通訊錄亦可撥者，書之於紙，藏之錢囊。
\end{enumerate}

\end{tcbpractice}

\booksection{場景一：手機遺失或被竊——黃金三十分鐘}

手機既失，首三十分鐘乃「黃金之時」，處之得宜，則財貨與隱私可保。依下序而行，一步毋越：

\textbf{失機後三十分鐘標準操作之流程}\par

\begin{enumerate}
  \item 借他人之機撥己之號（第 1 分鐘）：或止墜於近處，善人拾而接之，即可取回。
  \item 遠程鎖機（第 2-5 分鐘）：借他人之機或電腦登入搜尋之服務： iPhone：iCloud.com > 尋尋 > 擇爾之設備 > 標為遺失
  \item 安卓諸牌：華為 cloud.huawei.com / 小米 i.mi.com / OPPO cloud.oppo.com 等
  \item 掛失 SIM 卡（第 5-10 分鐘）：撥運營商客服之電話（移動 10086 / 聯通 10010 / 電信 10000），報手機已失，請停 SIM 卡之務。此步最要——騙者得爾之 SIM 卡，即可收驗證之碼，盡改爾之諸密。
  \item 凍結支付之帳（第 10-20 分鐘）： 支付寶：95188 掛失，或於他人之機登支付寶 > 我之 > 設定 > 安全中心 > 緊急凍結
  \item 微信：撥 110 報警後，聯微信客服，或於微信安全中心 > 凍結帳號
  \item 銀行 App：撥銀行卡背面之客服電話，請凍網銀
  \item 告親友（第 20-25 分鐘）：發於朋友圈，或託家人傳訊，告以手機已失，防騙者假爾之名以欺人。
  \item 報警（第 25-30 分鐘）：撥 110，或詣最近之派出所報案，取其回執。他日保險之賠、銀行之爭，皆需此執。
\end{enumerate}

\booksection{場景二：為人所騙——七十二小時黃金處理之期}

既覺受騙，要於鎮定。慌亂而誤決，必致再損。依下步而行：

\begin{enumerate}
  \item 立止一切轉帳之事：無論騙者何言，勿復轉一錢。「再轉一筆保金，前錢可返」者，二次行騙之常語也。
  \item 盡存其證：聊天之錄勿刪！通話之錄勿刪！截圖以存聊天之錄、轉帳之錄、對方帳號之訊。此皆報案追資之要證也。
  \item 急撥 96110：此全國反詐之專線。接線員將導爾緊急止付——若轉帳未久（常於二十四小時內），銀行或可凍對方之帳而截其資。
  \item 撥 110 報警：攜手機、身份證、轉帳之憑，詣最近之派出所錄供。
  \item 聯銀行、支付之臺：告以受騙之情，請凍涉案之帳。若信用卡為盜刷，請銀行啟拒付之程。
  \item 盡改諸密：於另一安全之器，改爾 Apple ID、品牌之帳、微信、支付寶、銀行 App 之密，並啟雙重認證。
\end{enumerate}

\booksection{場景三：手機染毒或木馬——奈何？}

若爾之機現此諸症——廣告無故自出、奇 App 自裝、話費暴增、流量異耗——則中惡意軟件之疑甚大：

\begin{itemize}
  \item 第一步：入安全模式——安卓之機長按電源鍵以關機，復長按以開機；開機既震，長按音量減鍵勿釋，至入安全模式而止（左下角現「安全模式」之字）。安全模式中，第三方 App 皆禁，便於排查。
  \item 第二步：卸可疑之 App——安全模式中，入設定 > 應用管理，以時間序之，得最近所裝陌生之 App 而卸之。
  \item 第三步：裝安全之 App 以掃之——下「騰訊手機管家」或「360手機衛士」，全盤一掃。
  \item 第四步：復原廠之設（終極之方）——若上法皆無效，先備要數據，乃復原廠。慎之：復原廠則諸數據盡刪，照片、通訊錄諸要檔，必先備之。
  \item iPhone 幾不染毒——iOS 沙盒之制，惡意軟件難存。若爾之 iPhone 有異，多為某 App 之 bug 或設定之失，非毒也。復原廠之設，可解九成九之患。
\end{itemize}

\booksection{場景四：帳號被竊——速復掌控之權}

\begin{enumerate}
  \item 急改其密：若尚能登其帳，第一時改密。新密須與舊密全異。
  \item 踢去他器：於帳號安全之設中得「登錄設備管理」，凡陌生之器，盡踢下線。
  \item 查登錄之記：覽近七日登錄之記，若有異地之登，則帳果被竊矣。
  \item 啟雙重認證：若未嘗啟，今即啟之。此防帳再竊之至效者也。
  \item 查所綁之訊：視帳號所綁之手機號、郵箱有無被改。若已改，以帳號申訴之功復之。
  \item 查支付之設：若被竊者為微信、支付寶，則視免密支付、自動扣款有無被篡，銀行卡有無見加於陌生之帳。
\end{enumerate}

\booksection{場景五：得勒索之訊——勿懼，勿付}

勒索之術有二：(1)「吾有爾不雅之影片、照片，不予錢則盡發爾通訊錄之人」；(2)「爾機已為黑客所制，不予錢則鎖死」。遇此：

\begin{itemize}
  \item 九成九皆詐也：彼實無爾絲毫之私。以同辭羣發數萬人，總有懼而付錢者。
  \item 勿復，勿付：一復則告彼曰「此號有人用」，爾遂被標為「可勒索之之」，其擾滋多。
  \item 截圖存證：遂拉黑其人。
  \item 改其密：若有憂，改爾主要帳號之密，以自安。
  \item 若果有洩私之虞：報警。警方辦此類之案，素有經驗。
\end{itemize}

\booksection{場景六：AI 換臉、AI 擬聲之詐——眼見不復為實}

自 2024 年，AI 換臉、擬聲之技，已能使騙者於視訊之中化為爾家人、長上、朋友之面與聲。「娘，兒出事矣，速轉錢來」——視中之人，貌與聲皆肖爾子，實則偽也。此當今至難防亦至危之詐術之一也。

\textbf{第一步：識其破綻——AI 偽造四常見之漏}\par

\begin{itemize}
  \item 眨眼之數異：真人每分鐘眨眼十五至二十，AI 換臉者多眨之過少，或其節僵。
  \item 口型與聲不諧：諦視口型與所聞之聲，有零點幾秒之差。
  \item 面際模糊：臉之輪廓、耳畔、髮際，易現「毛邊」或閃爍。
  \item 光影之向不侔：面上高光、陰影與背景照光之向不合，或側面時五官隨轉失其度。
\end{itemize}

\begin{tcbinfo}{第二步：三術驗其真偽}
\begin{itemize}
  \item 令作特定之動：使彼轉頭、掩面、以手障口而復開合。AI 換臉於障蔽與大角度轉頭之際最易露餡——面或變形，或追之不及。
  \item 問惟彼知之事：設一虛問，如「前日吾家之犬送何醫院視之？」——若果騙者，則不能對，或支吾其詞。
  \item 回撥已知之號：斷視訊，以爾通訊錄所存之號徑撥之。騙者不能制真之機號，一撥即敗。
\end{itemize}

\end{tcbinfo}

\begin{tcbwarn}{第三步：既覺受騙，報警之流程}
\begin{enumerate}
  \item 立止轉帳：雖急亦姑止，凡催爾「速轉帳、毋掛斷」者，殆皆騙也。
  \item 盡存其證：錄屏、截圖，存視訊之畫、聊天之錄、對方之帳號、轉帳之憑——此技術鑑定之要材也。
  \item 撥 96110 反詐專線：告以「AI 換臉、擬聲」之詐，接線員將啟緊急止付。
  \item 撥 110 報警：攜證詣派出所錄供，特陳其用 AI 偽造之技，以便警方立案定性。
  \item 告親友：騙者或以偽爾之視訊欺人，先於朋友圈、家族之羣發一提醒。
\end{enumerate}

\end{tcbwarn}

\booksection{場景七：SIM swap 之攻（SIM 劫持）——爾號為人所「竊」}

此攻人多未聞，而其害至巨。SIM swap 者（SIM 卡劫持），謂攻者偽造爾身份之訊，冒爾之名詣營業廳補辦新 SIM 卡。補辦既成，爾機之信號驟失，而攻者之機插彼新卡，盡收爾短信驗證之碼——支付寶、銀行、微信，門戶盡開矣。

\begin{tcbinfo}{何以覺之？SIM swap 三預警之兆}
\begin{itemize}
  \item 機忽示「無服務」、「無 SIM 卡」，而爾未嘗動卡，重啟亦無效——此多為被劫之頃。
  \item 得運營商「補卡既成」或「業務既辦」之短信——而爾未嘗辦也。
  \item 驗證之碼驟不得，他短信亦不得——以其短信被劫至攻者之卡也。
\end{itemize}

\end{tcbinfo}

\begin{tcbwarn}{防護之術——並鎖爾之 SIM 卡}
\begin{itemize}
  \item 設 SIM 卡之 PIN 碼：其徑 設定 > 電話 > SIM 卡 PIN 碼。既啟，每重啟或換卡，必輸 PIN 而後能用此卡。攻者雖得卡，亦無所用之。
  \item 啟運營商再驗之制：登中國移動、聯通、電信之 App，得「安全中心」或「業務辦理驗證」，啟非本人補卡之短信提醒與業務辦理再驗之密。如為他人補爾之卡，營業之員必求身驗之加焉。
  \item 要帳勿徒恃短信之驗：銀行、Apple ID 諸高值之帳，宜先用認證器 App（Google Authenticator、微軟 Authenticator）或硬體之鑰，毋用短信驗證之碼。
  \item 身份證影本加印水印：辦務所需之影本，明書「僅供某務之用」——攻者多賴洩漏之身份訊以冒爾。
\end{itemize}

\end{tcbwarn}

\textbf{被劫之後應急之處——唯速為要}\par

\begin{enumerate}
  \item 急撥運營商客服掛失：移動 10086 / 聯通 10010 / 電信 10000，請立停機並凍補卡之務，以阻攻者續用爾號。
  \item 以他器登要帳：用電腦或家人之機，急改銀行、支付寶、微信、Apple ID 之密，並踢陌生之器。
  \item 凍結支付之帳：支付寶 95188 掛失，微信安全中心凍之，銀行客服凍網銀。
  \item 持身份證詣營業廳：本人攜身份證至自營之廳補辦新卡，復其號之用。
  \item 報警：SIM swap 已構「侵犯公民個人信息之罪」，攜證（停機時間之記錄、資金之損失、補卡之記錄）以報。
  \item 遍察之：號既復，逐一視近期登錄之記與資金之流水，驗有無被冒用之跡。
\end{enumerate}

\begin{tcbwarn}{防勝於治——日用五事}
\begin{enumerate}
  \item 鎖屏之密與雙重認證，二者皆啟
  \item 「尋尋吾之機」務必啟
  \item SIM卡設PIN碼（設置 > 電話 > SIM卡PIN碼），丟了手機他人不能用汝卡
  \item 不裝來歷不明之 App，不點陌生之鏈
  \item 定期備份要數據（照片、通訊錄、聊天之錄）
\end{enumerate}

\end{tcbwarn}

\begin{tcbpractice}{來動手：製爾「應急聯絡之卡」}
\begin{enumerate}
  \item 取紙一張，書下列諸號：運營商客服（移動 10086 / 聯通 10010 / 電信 10000）、銀行客服、96110 反詐專線、110 報警之電話、家人之電話。
  \item 書爾 Apple ID、品牌之帳所註冊之郵箱，與備用聯絡之法。
  \item 置此紙於錢囊，或貼於冰箱——毋存於機，蓋事起之時，機或不於側也。
  \item 今即查之：「尋尋吾之機」已啟乎？SIM 卡 PIN 碼已設乎？
\end{enumerate}

\end{tcbpractice}

\begin{tcbsummary}{本章小結}
\begin{itemize}
  \item 失機黃金三十分鐘：借機撥號遠程鎖機掛失 SIM 卡凍結支付告親友報警
  \item 被騙七十二小時：止轉帳存其證撥 96110報警聯銀行改其密
  \item 機染毒：安全模式卸可疑 App安全之掃復原廠（終極之方）
  \item 帳被竊：改密踢器查登錄啟雙重認證查所綁查支付
  \item 勒索之訊九成九皆詐，勿復勿付，截圖而拉黑
  \item AI 換臉、擬聲之詐：觀眨眼、口型、邊際、光影之破綻，令彼轉頭掩面，問私密之事，回撥已知之號
  \item SIM swap 之劫：機忽「無服務」即警；設 SIM 卡 PIN、啟運營商再驗、要帳改用認證器 App
  \item 日用防患五事：鎖屏+2FA、尋尋手機、SIM 卡 PIN、不裝來歷不明之 App、定期備份
\end{itemize}

\end{tcbsummary}
